% Five replacements in Online Appendix C.6. The payoff oracle, the sweep design,
% identity and deduplication, reporting, and the output schemas stand, subject to
% the dictionary wording in the last replacement. No \NUM placeholders.
% In this subsection "pooling" as a solver branch (line 2784 and the branch names)
% means zero informed orders. That usage stays.

% REPLACES: online_appendix.tex lines 2700-2739 ("\textbf{Complete continuations.} ..." to "... fixed schedule.")
\textbf{Complete continuations.} A candidate at a reserve is accepted
only as a complete continuation in the sense of \eqref{eq:oa-oa-53}: the
investor's state- or class-contingent order distribution, the pricing
rule with every positive-mass price atom, the buyer's information at
each price (the atom posterior on an atom, the inverted posterior
elsewhere), the preparation rule by observed price and cost, and the
validation evidence. For a candidate order profile, compute its raw flow
posterior \(\mu_X\). When \(\Delta_T>0\), the benchmark inversion applies to
each price, because entry is at least \(\rho\) at every flow under the
declared cost laws. A positive-mass atom can still arise where \(\mu_X\) is
constant, as on a posterior plateau. Its posterior is the integral over the
full preimage of that price:

\[
\mu(P=P_a)=\frac{\int_{\{x:P(x)=P_a\}}a_H(x)\,dx}
{\int_{\{x:P(x)=P_a\}}[a_H(x)+a_L(x)]\,dx}
\tag{OA.79}\label{eq:oa-oa-79}
\]

for an atom at price \(P_a\). At a degenerate reserve \(P_a\) can be \(t_0\)
itself. The formula is used only when the denominator is positive.
The buyer's entry rule at that price uses this atom posterior and nothing finer; a
candidate whose buyer acts on the raw posterior inside an atom with one
posterior value loses nothing, and a candidate that assigns one price to flows
with distinct raw posteriors while the buyer acts on those posteriors is not
measurable with respect to the price and is rejected. The convenient
construction for independent costs,
\(P(\mu)=t_0+H_C(B(\mu))[t_L-t_0+\Delta_T\mu]\), prices every
flow by the inversion; it is strictly increasing wherever
\(\Delta_T>0\) and the cost law gives \(H_C(B(\mu))\ge\rho\).
It is the one pricing family the C.6 node solver
constructs, and the row says so: \texttt{pricing\_family\_coverage}
records that no other pricing family was scanned at that node and
\texttt{exhaustive\_pricing\_search} is false. That record is now a coverage
statement about order profiles and degenerate reserves, not about pooling
regions: \hyperref[oa-a-pools]{Section A.11} shows that at fixed orders with
\(\Delta_T>0\) no second price mapping survives the buyer check.
This does not establish uniqueness of the order profile at a reserve. Every
accepted candidate must pass the price-atom validation of C.0: atom
posterior from the full preimage, buyer optimality at the atom
posterior and pointwise on other prices, competitive pricing
pointwise in flow on both regions, no collision between an atom and the
rest of the price range, and every unilateral deviation against the whole
fixed schedule.

% REPLACES: online_appendix.tex lines 2741-2747 ("At reserves with \Delta_T=0, ..." to "... entry conditions.")
At reserves with \(\Delta_T=0\), no sale possible, or a degenerate
payoff, analyze the constant-price candidate directly. Zero orders use the
actual prior entry probability \(e_0=H_C(B(1/2))\), and their correctly
signed deviation coefficient is \(e_0\Delta_T/2-k\); under the declared
atomic law \(e_0\ge\rho\), so that coefficient starts from
\(\rho\Delta_T/2-k\), while a cost law with no atom at zero can still give a
smaller \(e_0\) at a reserve where gross profit is zero. When \(e_0=0\), no trade and no
entry are supported by constant pricing. All accepted candidates must
satisfy the original conditional pricing and entry conditions.

% REPLACES: online_appendix.tex lines 2749-2754 ("\textbf{Specified comparisons.} ..." to "... not a global reserve optimum.")
\textbf{Specified comparisons.} At the declared reserves, verify the
relevant strict no-trade or full-order bounds using the class payoff
spread and the entry value at the lowest belief \(e(m)=H_C(B_r(m))\),
which is at least \(\rho\) under the declared laws. Compute entry and
seller revenue and compare each alternative reserve with the original
one at the same strength. This reproduces the profitable alternatives,
not a global reserve optimum.

% REPLACES: online_appendix.tex lines 2758-2777 ("\textbf{Event grid.} ..." to "... record does not assume it.")
\textbf{Event grid.} The sweep and the event exercise share one grid
rule. Beyond the declared reserve mesh, every node set contains the
support and participation events \(0\), \(\ell\), \(r\), \(h\); in the
class economy \(\ell\pm\varepsilon_V\) and \(h\pm\varepsilon_V\); the
declared reserve alternatives; and the two exact events of the paper
appendix, the endpoint equality \(p_L=h\) and the
expensive-entry ceiling equality \(p_H=\sqrt{2r(h-c_H/M)-r^2}\), each
attached to the strength at which its defining regime applies. Every
event carries one-sided offsets \texttt{1e-4}, \texttt{1e-6}, and
\texttt{1e-8} on both sides where the offset stays in \([0,h]\) (or
\([0,h+\varepsilon_V]\)), so that the regime classification on each side
is tested separately. At \(p_L=h\) the offsets run only below the endpoint.
An exact event is stored by its defining relation
and recomputed from it at fifty digits; its decimal expansion is
reference only and never decides a tie. The tie rule is applied to the
identity: at \(p_L\) the zero-cost realization enters at every belief under
\(B(\mu)=0\), and at \(p_H\) the high-cost type enters on the upper
plateau under \(\tau=M\) with \(x^*=1\). A node whose inequality sign
cannot be resolved at that precision is written as unresolved. The same
continuation need not be supported on both sides of the interior event, and the
record does not assume it.

% REPLACES: online_appendix.tex lines 2900-2921 ("\texttt{low_cost_floor_relation} and ..." to "... not written as a no-entry pool.")
\texttt{low\_cost\_floor\_relation} and \texttt{ceiling\_relation} take
the values strict positive, strict negative, equality event, or
unresolved, from the fifty-digit evaluation. The first relation now compares
gross profit at the lowest belief, \(B_r(m)\) in the node's regime, with the
zero preparation cost; its equality event inside the reserve domain is the
endpoint \(p_L=h\). \texttt{candidate\_id} is
the attempt identity and \texttt{continuation\_id} the economic identity
defined above; \texttt{duplicate\_of} is empty or the representative's
candidate identity; \texttt{identity\_note} records
same-orders-different-atoms cases. \texttt{parameter\_set\_id} is the
canonical string of exact declarations (not a hash);
\texttt{institution\_id}, \texttt{information\_structure\_id}, and
\texttt{tie\_rule\_id} are the literal definitions of the sale rule, of
who observes what, and of the tie conventions. \texttt{order\_rule\_id}
is the canonical order profile; \texttt{price\_rule\_id} names the
pricing family, its price atoms, and the positive-entry price map.
\texttt{pricing\_family\_coverage} and
\texttt{exhaustive\_pricing\_search} are as declared in C.0.
\texttt{pool\_id}, \texttt{pool\_set\_definition},
\texttt{pool\_cutoff\_exact}, \texttt{price\_atom\_value},
\texttt{state\_H\_pool\_mass}, \texttt{state\_L\_pool\_mass},
\texttt{pool\_probability}, and \texttt{pool\_posterior} describe a
constant-price atom when one exists and are \texttt{n/a} otherwise. Under the
declared cost laws such an atom has entry at least \(\rho\) at each reserve
with \(\Delta_T>0\); a no-entry atom is possible only at a degenerate reserve.
A plateau atom in a flat posterior tail is recorded through the same columns,
with its preimage stated in \texttt{pool\_set\_definition}.
